\title{Controlling Text-to-Image Diffusion by Orthogonal Finetuning}

\begin{abstract}
Recent advances in large-scale text-to-image diffusion models have demonstrated remarkable proficiency in synthesizing lifelike images based on textual descriptions. However, efficiently steering these sophisticated models for a variety of downstream tasks presents a significant open question. In response, we propose Orthogonal Finetuning (OFT), a theoretically grounded adaptation technique for refining text-to-image diffusion models to suit specific downstream applications. Distinct from conventional strategies, OFT can formally guarantee the preservation of hyperspherical energy, a measure representing the relational structure among neurons situated on the unit hypersphere. We establish that this preservation is key to maintaining the semantic generation capacity of text-to-image diffusion models. Furthermore, to promote stable finetuning, we introduce Constrained Orthogonal Finetuning (COFT), which enforces a supplementary radius constraint on the hypersphere. Our work focuses on two representative finetuning scenarios: subject-driven generation, which entails synthesizing subject-specific images with a handful of reference photos and a prompt; and controllable generation, which involves incorporating auxiliary control cues into the generation process. Through empirical evaluation, we demonstrate that the OFT approach surpasses existing alternatives in both the quality of generated images and the speed of convergence.
\end{abstract}

\section{Introduction}

Text-to-image diffusion models~\cite{saharia2022photorealistic,ramesh2022hierarchical,rombach2022high} have set new benchmarks in the generation of high-fidelity images guided by textual instructions. Yet, even with strong baseline capabilities, textual input alone often falls short in delivering the fine-level and precise controls required for certain generative outcomes. This work seeks to address these challenges by exploring two core categories of text-to-image generation tasks:

\begin{itemize}[leftmargin=*,nosep]

    \item \textbf{Subject-driven generation}~\cite{ruiz2023dreambooth}: Provided with a limited collection of images depicting a specific subject, this task requires creating novel images of that subject situated in new contexts under the direction of a text prompt.
    \item \textbf{Controllable generation}~\cite{zhang2023adding,mou2023t2i}: When additional control information (\emph{e.g.}, canny edge maps, segmentation masks) is given, the objective is to generate images that comply both with this guidance and a given text description.
\end{itemize}

At their core, both tasks are concerned with strategies for finetuning text-to-image diffusion models in a way that maintains the quality and diversity established during pretraining. The primary objectives for a robust finetuning approach are: (1) \emph{training efficiency}, which entails minimizing both the number of trainable parameters and required epochs, and (2) \emph{generalizability preservation}, ensuring the model continues to generate high-fidelity and varied outputs. To address these goals, common finetuning methods typically involve either making incremental adjustments to neuron weights using a modest learning rate (\emph{e.g.}, \cite{ruiz2023dreambooth}), or appending a lightweight module composed of re-parameterized neuron weights (\emph{e.g.}, \cite{hulora2022,zhang2023adding}). Although straightforward in their application, these techniques do not provide assurances that the original generative capabilities from pretraining are retained. Additionally, there is a noticeable gap in foundational insights for constructing optimal finetuning procedures and determining effective hyperparameter choices, such as the number of epochs required. A principal challenge is the absence of a concrete metric that captures the extent to which the finetuned model preserves pretrained generative qualities. Currently, prevalent methods rely on the implicit premise that minimizing the Euclidean distance between the parameters of the finetuned and pretrained models will yield better retention of generative performance, which motivates the use of low learning rates during finetuning. However, while this assumption may sometimes be valid, we contend that Euclidean proximity alone does not sufficiently account for the semantic integrity of the model. As such, adopting a more structured metric to evaluate the relationship between the finetuned and pretrained models could significantly enhance both the preservation of pretrained generative features and the overall stability during finetuning.

Drawing inspiration from empirical findings that hyperspherical similarity effectively captures semantic relationships~\cite{liu2017deep,liu2018decoupled,chen2020angular}, our approach employs hyperspherical energy~\cite{liu2018learning} to quantify the relational dynamics between neurons within the same layer. Specifically, hyperspherical energy is calculated as the aggregated hyperspherical similarities (\emph{e.g.}, cosine similarities) among all pairs of neurons, thus reflecting the extent of neuron dispersion over the unit hypersphere~\cite{liu2021learning}. Our premise is that optimal finetuning should aim to minimize any change in hyperspherical energy relative to the original, pretrained model. A straightforward strategy would be to introduce a regularization term to keep hyperspherical energy constant throughout finetuning, though this method does not ensure that the change in hyperspherical energy is sufficiently controlled. To address this, we leverage the invariance property of hyperspherical energy: when the same orthogonal transformation is applied to all neurons, their pairwise hyperspherical similarity remains intact. Building on this property, we introduce Orthogonal Finetuning~(OFT), which enables the adaptation of large text-to-image diffusion models for downstream applications while maintaining their hyperspherical energy. The core of this technique is to learn a single orthogonal transformation shared across the layer, thereby preserving all pairwise angular relationships among neurons. Conceptually, OFT can be interpreted as redefining the canonical coordinate axes for the neurons in a given layer. Furthermore, taking into account the established relationship between smaller Euclidean distances to the pretrained model and better retention of pretraining performance, we advance an extension—Constrained Orthogonal Finetuning~(COFT)—which confines the finetuned model to remain within a hypersphere of fixed radius around the pretrained neurons.

The rationale underlying the effectiveness of orthogonal transformation in neuron finetuning draws, in part, from the properties of the 2D Fourier transform, which decomposes an image into its magnitude and phase spectra. Critically, the phase spectrum—capturing angular relations between the input and basis functions—encapsulates most of the semantic content. For instance, reconstructing an image using its original phase spectrum paired with an arbitrary magnitude spectrum retains the key semantic features of that image. This observation indicates that modifying the orientation of neurons is essential for introducing semantic variations in generated images, which is central to both subject-driven and controllable image synthesis. Nevertheless, unrestricted modification of neuron orientations risks significantly impairing the generative capacity acquired during pretraining. To address this, we advocate for maintaining the angles between all pairs of neurons, premised on the conjecture that such angular relationships encode the neural network's learned knowledge. Accordingly, we employ layer-wise orthogonal transformations that preserve inter-neuron angles, ensuring the hyperspherical energy remains unchanged across layers.

Additionally, we are informed by orthogonal over-parameterized training~\cite{liu2021orthogonal}, which entails orthogonally transforming the weights of a randomly initialized network for classification tasks. The appeal of this approach stems from the fact that random initialization tends to yield neural networks with low expected hyperspherical energy, a property \cite{liu2021orthogonal} leverages by enforcing the preservation of this low energy during training—leading to improved generalization performance~\cite{liu2018learning,lin2020regularizing}. Their findings demonstrate that orthogonal transformations afford sufficient flexibility to produce classification networks capable of generalizing effectively. In contrast, our work centers on finetuning text-to-image diffusion models to enhance both controllability and generative prowess in downstream tasks. OFT differs from~\cite{liu2021orthogonal} along two primary dimensions. First, rather than targeting hyperspherical energy minimization, OFT aims to retain the pretrained model’s hyperspherical energy, thereby safeguarding its intrinsic representational structure throughout the finetuning process—since minimizing hyperspherical energy may compromise semantic integrity in diffusion models. Second, OFT incorporates a constraint to minimize divergence from the pretrained model, resulting in a more restricted version, unlike the unconstrained approach of~\cite{liu2021orthogonal}. Effective finetuning hinges on striking an optimal balance between adaptability and stability; we contend that the OFT framework succeeds in achieving this equilibrium. Our principal contributions are enumerated below:

\begin{itemize}[leftmargin=*,nosep]

    \item We introduce a new finetuning technique, Orthogonal Finetuning (OFT), aimed at enhancing the controllability of text-to-image diffusion models. Additionally, to bolster stability, we develop a constrained variant that restricts the angular change relative to the original pretrained parameters.
    \item In contrast to prevalent finetuning strategies, OFT adjusts model parameters while mathematically ensuring the preservation of hyperspherical energy—a property we empirically identify as vital for maintaining the generative semantics encoded by the pretrained network.
    \item We evaluate OFT in two specific scenarios: subject-driven generation and controllable generation. Through extensive experimental analysis, we demonstrate that OFT significantly outperforms prior approaches with respect to output quality, convergence rate, and stability during finetuning. Notably, OFT is also more sample-efficient, attaining strong results with fewer training examples and epochs.
    \item To assess controllable generation, we propose a novel control consistency metric. This metric involves extracting the control signal from the produced images and comparing it to the original control input. Our results using this metric underscore OFT's robust controllability.
\end{itemize}

\section{Related Work}

\textbf{Text-to-image diffusion models}. Recent advancements~\cite{saharia2022photorealistic,ramesh2022hierarchical,rombach2022high,gu2022vector,nichol2022glide} have significantly elevated text-to-image synthesis, largely attributable to progress in diffusion generative modeling~\cite{ho2020denoising,song2021score,song2021denoising,dhariwal2021diffusion} alongside the maturation of vision-language representation frameworks~\cite{su2020vlbert,lu2019vilbert,radford2021learning,wang2021simvlm,li2022blip,li2023blip,alayrac2022flamingo,schuhmann2022laion}. For instance, GLIDE~\cite{nichol2022glide} and Imagen~\cite{saharia2022photorealistic} operate diffusion processes directly within the image (pixel) domain; GLIDE jointly learns a text encoder together with a diffusion prior using paired data, whereas Imagen employs a frozen pretrained text encoder. Models like Stable Diffusion~\cite{rombach2022high} and DALL-E2~\cite{ramesh2022hierarchical} shift the diffusion to the latent domain: Stable Diffusion uses VQ-GAN~\cite{esser2021taming} to learn a visual codebook serving as its latent space, while DALL-E2 utilizes CLIP~\cite{radford2021learning} to produce a unified latent space for both images and text. Apart from diffusion architectures, generative adversarial networks~\cite{reed2016generative,zhang2017stackgan,xu2018attngan,li2019controllable} and autoregressive models~\cite{ramesh2021zero,ding2021cogview,wu2022nuwa,yuscaling} have also been widely adopted in the text-to-image field. OFT is a model-agnostic finetuning method, compatible with any text-to-image diffusion framework.

\textbf{Subject-driven generation}. To avoid unwanted alterations of subject properties, works such as~\cite{avrahami2022blended,nichol2022glide} introduce user-provided masks as additional constraints. Inversion techniques~\cite{choi2021ilvr,dhariwal2021diffusion,ramesh2022hierarchical,gal2022image} enable contextual changes while retaining subject identity. Meanwhile,~\cite{hertz2022prompt} supports both localized and global editing with no requirement for input masks. Nevertheless, most existing solutions are limited in faithfully preserving identity-specific traits. Pivotal Tuning~\cite{roich2022pivotal} enhances identity consistency by optimizing around an inverted latent vector using regularization, while~\cite{nitzan2022mystyle} formulates a bespoke generative face prior from a set of facial images. The instance generation method in~\cite{casanova2021instance} creates various appearances for a given entity, though it sometimes fails to retain detailed instance characteristics. DreamBooth~\cite{ruiz2023dreambooth} adapts a text-to-image diffusion model to specific subjects by incorporating a custom token and using few example images; finetuning is performed via a reconstruction loss and a class-driven prior preservation objective. OFT is developed on top of the DreamBooth protocol, substituting conventional fine adjustment via low learning rates with finetuning through orthogonal parameter transformations.

\textbf{Controllable generation}. Image-to-image translation is a classic instance of controllable generation, with the majority of earlier approaches utilizing conditional generative adversarial networks~\cite{isola2017image,zhu2017toward,wang2018high,park2019semantic,choi2018stargan}. Recently, diffusion models have also gained traction in image-to-image translation tasks~\cite{saharia2022palette,voynov2022sketch,wang2022pretraining}. More recent advancements include ControlNet~\cite{zhang2023adding}, which refines a pretrained diffusion model through finetuning and adjustment for additional control cues, resulting in state-of-the-art controllable generation. Similarly, T2I-Adapter~\cite{mou2023t2i}, a parallel and related study, achieves enhanced control in generated images by adapting a pretrained diffusion model through finetuning. Adhering to the experimental protocols of \cite{zhang2023adding,mou2023t2i}, we leverage OFT to adapt pretrained diffusion models, demonstrating improved controllability while requiring less training data and fewer finetuned parameters. Importantly, OFT introduces no extra computational demands at inference time.

\textbf{Model finetuning}. The practice of refining large pretrained models for downstream applications is becoming standard~\cite{devlin2018bert,brown2020language,he2022masked}. Within the finetuning paradigm, a variety of adaptation techniques (\emph{e.g.}, \cite{hulora2022,houlsby2019parameter,pfeiffer2020adapterfusion}) have been thoroughly explored in NLP research. LoRA~\cite{hulora2022} is particularly closely related to OFT, as it models weight adaptation during finetuning with a low-rank additive update. In contrast, OFT implements a multiplicative, layer-consistent orthogonal transformation to update neuron weights, with a theoretical guarantee that intra-layer neuron pairwise angles are preserved, thereby enhancing training stability.

\section{Orthogonal Finetuning}

\subsection{Why Does Orthogonal Transformation Make Sense?}

To motivate the use of orthogonal transformations in finetuning text-to-image diffusion models, we first consider two foundational questions: (1) the rationale for finetuning the directional component (i.e., angle) of neurons, and (2) the suitability of orthogonal transformations for modifying those directions.

To address the first question, we are motivated by empirical findings from \cite{liu2018decoupled,chen2020angular}, which indicate that differences in angular features provide a strong representation of the semantic gap. SphereNet~\cite{liu2017deep} demonstrates that constraining all neurons of a neural network to lie on the unit hypersphere can preserve, and sometimes even enhance, the network’s expressive capacity and generalization ability, suggesting that the directional component of neurons is crucial for encapsulating key data properties. To illustrate the role of neuron angles, we carry out a toy experiment (see Figure~\ref{fig:toy_ae}) where a conventional convolutional autoencoder is trained on a dataset of flower images. During training, standard inner products are used to compute the feature map ($z$ stands for the output element obtained via convolving kernel $\bm{w}$ over input $\bm{x}$ within the sliding window). For evaluation, we consider three distinct strategies to reconstruct the feature map: (a) employing the original inner product as in training, (b) retaining only the magnitude, and (c) extracting only the angular component. The outcomes presented in Figure~\ref{fig:toy_ae} reveal that neuron angles alone are nearly sufficient to reconstruct the input images, whereas magnitude conveys negligible semantic content. It is important to note that the cosine-based (c) activation is not used during the learning phase, which exclusively relies on the inner product. These findings suggest that the semantic essence of the input images predominantly resides in the directions of neuron weights. Consequently, adjusting neuron orientations offers a highly effective means to alter semantic information.

Regarding the second question, a straightforward method for refining neuron directions involves collectively rotating or reflecting all neurons within the same layer, which inherently corresponds to an orthogonal transformation. While one could pursue more intricate angular transformations, such as rotating individual neurons by distinct angles, we observe that orthogonal transformations strike an advantageous compromise between adaptability and structure. Additionally, prior work \cite{liu2021orthogonal} has established that orthogonal transformations offer ample expressive power for training neural networks. To substantiate this claim, we present an experiment designed to examine the regularization effect provided by orthogonal constraints. Employing a controllable generation protocol in line with ControlNet~\cite{zhang2023adding}, we showcase the outcomes in Figure~\ref{fig:ortho}. Our results indicate that when the standard OFT includes the orthogonality constraint, stable and precise control is achieved after training for 20 epochs. By contrast, OFT lacking this constraint is unable to generate plausible images or exhibit meaningful control, thereby underscoring the critical function of orthogonality in OFT.

\subsection{General Framework}

Standard finetuning approaches generally rely on gradient descent with a reduced learning rate to update model parameters or selected layers. This small learning rate acts as an implicit regularizer, constraining parameter updates to remain close to the original pretrained weights, effectively optimizing $\thickmuskip=2mu \medmuskip=2mu \| \bm{M} - \bm{M}^0\|$, where $\bm{M}$ and $\bm{M}^0$ denote finetuned and pretrained weights, respectively. Although this form of regularization appears reasonable, it may not provide sufficient control when adapting large-scale models. To enhance this, LoRA imposes an explicit low-rank structure on the update, specifically enforcing $\thickmuskip=2mu \medmuskip=2mu \text{rank}(\bm{M} - \bm{M}^0)=r'$ with $r'$ set to a small value. In contrast, OFT adds a different restriction, targeting the pairwise similarity among neurons, such that $\thickmuskip=2mu \medmuskip=2mu \| \text{HE}(\bm{M})-\text{HE}(\bm{M}^0) \|=0$, where $\text{HE}(\cdot)$ computes the hyperspherical energy for a given weight matrix.

To elucidate this, take as example a fully connected layer $\thickmuskip=2mu \medmuskip=2mu \bm{W}=\{\bm{w}_1,\cdots,\bm{w}_n\}\in\mathbb{R}^{d\times n}$, where each neuron $\bm{w}_i\in\mathbb{R}^{d}$ and $\bm{W}^0$ corresponds to the pretrained parameters. The output for this layer, $\thickmuskip=2mu \medmuskip=2mu \bm{z}\in\mathbb{R}^n$, is given by $\thickmuskip=2mu \medmuskip=2mu \bm{z}=\bm{W}^\top\bm{x}$ for an input $\thickmuskip=2mu \medmuskip=2mu \bm{x}\in\mathbb{R}^d$. In this context, OFT can be framed as minimizing the change in hyperspherical energy between the updated and original weights:

\begin{equation}\label{eq:oft_principle}
\footnotesize
\min_{\bm{W}} \left\| \text{HE}(\bm{W})-\text{HE}(\bm{W}^0)\right\|~~\Leftrightarrow~~\min_{\bm{W}} \bigg{\|} \sum_{i\neq j} \|\hat{\bm{w}}_i-\hat{\bm{w}}_j\|^{-1}-\sum_{i\neq j} \|\hat{\bm{w}}^0_i-\hat{\bm{w}}^0_j\|^{-1}\bigg{\|}
\end{equation}

Here, the notation $\thickmuskip=2mu \medmuskip=2mu \hat{\bm{w}}_i:=\bm{w}_i/\|\bm{w}_i\|$ denotes the normalized $i$-th neuron, and the hyperspherical energy for $\bm{W}$ is calculated as $\thickmuskip=2mu \medmuskip=2mu \text{HE}(\bm{W}):= \sum_{i\neq j} \|\hat{\bm{w}}_i-\hat{\bm{w}}_j\|^{-1}$. It is evident that the minimum value for Eq.~\eqref{eq:oft_principle} is zero, which occurs if $\bm{W}$ can be derived from $\bm{W}^0$ via an orthogonal transformation, namely, $\thickmuskip=2mu \medmuskip=2mu \bm{W}=\bm{R}\bm{W}^0$, with $\thickmuskip=2mu \medmuskip=2mu \bm{R}\in\mathbb{R}^{d\times d}$ as an orthogonal matrix (a determinant of $1$ represents rotation, $-1$ indicates reflection). This central premise underpins OFT: rather than tuning individual parameters freely, finetuning is accomplished by learning a set of orthogonal matrices—one for each layer—to modify all neurons collectively. In practical terms, OFT formulates the learning of an orthogonal

\begin{equation}\label{eq:oft_general}
    \footnotesize
    \bm{z} = \bm{W}^\top\bm{x} = (\bm{R}\cdot\bm{W}^0)^\top\bm{x},~~~\text{s.t.}~\bm{R}^\top\bm{R} = \bm{R}\bm{R}^\top = \bm{I}
\end{equation}

In this formulation, $\bm{W}$ indicates the weight matrix after OFT fine-tuning, while $\bm{I}$ refers to the identity matrix. Figure~\ref{fig:block} visually summarizes the OFT procedure. In alignment with the approach of zero initialization as applied in LoRA, OFT must be configured such that fine-tuning initiates exactly from the pretrained parameters $\bm{W}^0$. To enforce this, we set the initial value of the orthogonal matrix $\bm{R}$ to the identity matrix, which guarantees that the model's starting parameters match the original pretrained state. Maintaining the orthogonality of $\bm{R}$ can be achieved by leveraging differentiable orthogonalization techniques introduced by \cite{liu2021orthogonal,lezcano2019cheap}. The following discussion explores practical strategies to preserve orthogonality efficiently.

\subsection{Efficient Orthogonal Parameterization}\label{sect:eff_ortho}

Conventional orthogonalization procedures, such as the Gram-Schmidt process, are differentiable but incur substantial computational costs in real-world scenarios~\cite{liu2021orthogonal}. To improve computational efficiency, we opt for Cayley parameterization to construct the orthogonal matrix. Concretely, the orthogonal matrix is formed as $\thickmuskip=2mu \medmuskip=2mu \bm{R} = (\bm{I}+\bm{Q})(I-\bm{Q})^{-1}$, where the matrix $\bm{Q}$ is skew-symmetric, that is, $\thickmuskip=2mu \medmuskip=2mu \bm{Q} = -\bm{Q}^\top$. While this method is efficient, it imposes the restriction that only orthogonal matrices with determinant $1$—members of the special orthogonal group—can be represented. Nevertheless, in our experiments, this constraint does not adversely impact empirical performance. Although Cayley transforms enable differentiable parameterization of the orthogonal matrix, as the dimension $d$ increases, representing $\bm{R}$ in this way remains parameter-intensive. To alleviate this, we propose to parameterize $\bm{R}$ as a block-diagonal matrix consisting of $r$ blocks, as described below:

\begin{equation}
    \footnotesize
    \bm{R} = \text{diag}(\bm{R}_1, \bm{R}_2, \cdots, \bm{R}_r) = \begin{bmatrix}
    \bm{R}_1\in \text{O}(\frac{d}{r}) &  &\\
    &\ddots & \\
    & & \bm{R}_r\in \text{O}(\frac{d}{r})
    \end{bmatrix}\in \text{O}(d)
\end{equation}

Here, $\text{O}(d)$ represents the orthogonal group in $d$ dimensions. The matrices $\thickmuskip=2mu \medmuskip=2mu \bm{R} \in \mathbb{R}^{d \times d}$ and, for each $i$, $\thickmuskip=2mu \medmuskip=2mu \bm{R}_i \in \mathbb{R}^{d/r \times d/r}$ are orthogonal. In the special case when $\thickmuskip=2mu \medmuskip=2mu r = 1$, the block-diagonal orthogonal matrix reduces to the standard, unconstrained orthogonal case. For a $d \times d$ orthogonal matrix, the total number of free parameters is $\thickmuskip=2mu \medmuskip=2mu d(d-1)/2$, yielding computational complexity of $\mathcal{O}(d^2)$. When using a block-diagonal orthogonal matrix with $r$ blocks, the parameter count becomes $\thickmuskip=2mu \medmuskip=2mu d(d/r-1)/2$, or complexity $\mathcal{O}(d^2/r)$. Moreover, by reusing the same block matrix for all blocks, i.e., by enforcing $\thickmuskip=2mu \medmuskip=2mu \bm{R}_i = \bm{R}_j$ for all $i \neq j$, parameter requirements decrease further to $\mathcal{O}(d^2/r^2)$. Importantly, all these methods still maintain the overall orthogonality of $\bm{R}$, thus the underlying hyperspherical structure is preserved without compromise.

To compare the parameter efficiency between OFT and LoRA, consider LoRA’s trainable parameter count, which with rank $r'$ amounts to $\thickmuskip=2mu \medmuskip=2mu r'(d+n)$. Suppose both $r$ and $r'$ scale with the width $d$ (for example, $\thickmuskip=2mu \medmuskip=2mu r=r'=\alpha d$ with $\thickmuskip=2mu \medmuskip=2mu 0<\alpha\leq 1$ for some $\alpha$), then LoRA’s parameter complexity becomes $\mathcal{O}(d^2+dn)$, whereas OFT’s becomes $\mathcal{O}(d)$. To clarify, consider a matrix of size $\thickmuskip=2mu \medmuskip=2mu 128 \times 128$; with $\thickmuskip=2mu \medmuskip=2mu r'=8$, LoRA uses $2,048$ parameters, whereas OFT with $\thickmuskip=2mu \medmuskip=2mu r=8$ (without shared blocks) uses only $960$ parameters.

\subsection{Constrained Orthogonal Finetuning}

To further restrict the adaptation made during OFT, one can constrain the finetuned weights to remain within a narrow region surrounding the pretrained parameters. In COFT, this is achieved by computing activations as follows:

\begin{equation}\label{eq:coft_general}
    \footnotesize
    \bm{z}=\bm{W}^\top\bm{x}=(\bm{R}\cdot\bm{W}^0)^\top\bm{x},~~~\text{s.t.}~\bm{R}^\top\bm{R}=\bm{R}\bm{R}^\top=\bm{I},~\norm{\bm{R}-\bm{I}}\leq \epsilon
\end{equation}

This approach enforces both orthogonality of $\bm{R}$ and restricts its deviation from the identity by at most $\epsilon$. The Cayley parameterization discussed in Section~\ref{sect:eff_ortho} can be utilized to ensure orthogonality. However, explicitly imposing the $\epsilon$-ball constraint on the Cayley-parameterized $\bm{R}$ is challenging. To better understand the Cayley transform’s properties, one may employ a Neumann series expansion to approximate $\thickmuskip=2mu \medmuskip=2mu \bm{R}=(\bm{I}+\bm{Q})(I-\bm{Q})^{-1}$ as $\thickmuskip=2mu \medmuskip=2mu \bm{R}\approx \bm{I} + 2\bm{Q} + \mathcal

Since the series converges under the operator norm, it is valid to shift the constraint $\thickmuskip=2mu \medmuskip=2mu\|\bm{R}-\bm{I}\|\leq\epsilon$ directly into the Cayley transform. Under this setting, the condition becomes $\thickmuskip=2mu \medmuskip=2mu \|\bm{Q}-\bm{0}\|\leq\epsilon'$, where $\bm{0}$ stands for the zero matrix and $\epsilon'$ is a distinct error threshold compared to $\epsilon$. This reformulated constraint on $\bm{Q}$ can be efficiently managed using projected gradient descent. For initialization that ensures $\bm{R}$ is the identity matrix, one simply sets $\bm{Q}$ to all zeros. Conceptually, COFT incorporates two explicit regularization mechanisms: constraining both hyperspherical energy changes and deviations from the original, pretrained parameters. While the latter is generally employed in many standard finetuning routines implicitly, COFT formalizes this as an explicit constraint. Empirically, although OFT achieves strong performance, we find that COFT provides enhanced finetuning robustness due to its overt handling of model deviation. Figure~\ref{fig:eps_coft} illustrates how different choices of $\epsilon$ influence COFT outcomes. Notably, as $\epsilon$ increases, the behavior of COFT converges to that of OFT, while smaller $\epsilon$ values make COFT increasingly align with the original pretrained text-to-image diffusion model.

\subsection{Re-scaled Orthogonal Finetuning}

We introduce a straightforward modification of OFT that additionally involves learning a scaling parameter for each neuron. The rationale stems from the observation that scaling neuron outputs does not impact their hyperspherical energy (as normalization removes magnitude effects). Concretely, the forward computation is given by $\thickmuskip=2mu \medmuskip=2mu\bm{z}=(\bm{R}\bm{W}^0\bm{D})^\top\bm{x}$\footnote{\textbf{Errata}: The NeurIPS camera-ready version incorrectly reported the re-scaled OFT forward formula as $\thickmuskip=2mu \medmuskip=2mu\bm{z}=(\bm{D}\bm{R}\bm{W}^0)^\top\bm{x}$. The original implementation, however, uses the correct expression and thus all results in Appendix~\ref{app:rescaled} remain valid.}, where $\thickmuskip=2mu \medmuskip=2mu \bm{D}=\text{diag}(s_1,\cdots,s_n)\in\mathbb{R}^{n\times n}$ is a diagonal matrix with each $s_i > 0$, and its entries are learned. This setup differs from the original OFT, where only $\bm{R}$ is optimized; here, both the diagonal matrix $\bm{D}$ and the orthogonal matrix $\bm{R}$ are trainable parameters. This re-scaled version extends OFT’s flexibility at the cost of a minimal increase in parameters. For our experiments, we retain the original OFT to isolate the benefits of orthogonal transformations alone, though supplementary evidence in Appendix~\ref{app:rescaled} shows that re-scaled OFT typically yields further gains.

\section{Intriguing Insights and Discussions}

\textbf{OFT’s model-agnostic nature}. The OFT technique is, in principle, broadly applicable to various neural network architectures. Although LoRA is often specialized to attention weight matrices in Transformers~\cite{hulora2022}, we restrict OFT's application to attention layers to facilitate fair benchmarking against LoRA in our evaluations. Beyond dense (fully connected) layers, OFT is inherently compatible with convolutional layers. The block-diagonal structure of $\bm{R}$ assumes particularly meaningful roles in convolutional settings—a property not shared by LoRA. By matching the block count in $\bm{R}$ with the number of input channels, each block modifies an individual channel's activations, analogous to tuning depth-wise convolution kernels~\cite{chollet2017xception}. Alternatively, with shared blocks, $\bm{R}$ executes an orthogonal transformation uniformly across all channels. We offer preliminary results demonstrating OFT’s suitability for convolutional layer finetuning in Appendix~\ref{app:convolution}.

\textbf{Connection to LoRA}. LoRA achieves preservation of pretrained information by incorporating a low-rank matrix, which restricts drastic alterations to the pretrained weight matrix. Conversely, OFT utilizes an orthogonal (full-rank) transformation applied to the pretrained weight matrix, ensuring that this operation does not compromise the integrity of pretraining knowledge. The forward computation for OFT may be reformulated as $\thickmuskip=2mu \medmuskip=2mu \bm{z}=(\bm{R}\bm{W}^0)^\top\bm{x}=(\bm{W}^0+(\bm{R}-\bm{I})\bm{W}^0)^\top\bm{x}$, where $\thickmuskip=2mu \medmuskip=2mu (\bm{R}-\bm{I})\bm{W}^0$ serves a similar role to LoRA's low-rank adaptation. When $\bm{W}^0$ possesses full rank, OFT can emulate low-rank adaptation in cases where $\thickmuskip=2mu \medmuskip=2mu \bm{R}-\bm{I}$ remains low-rank. In parallel with LoRA’s use of the rank parameter $r'$, OFT employs a diagonal block size $r$ to decrease the number of parameters requiring training. Importantly, LoRA and OFT embody distinct strategies for achieving parameter efficiency: LoRA leverages low-rank representations, whereas OFT focuses on exploiting sparsity patterns (specifically, block-diagonal orthogonal structures).

\textbf{Why OFT converges faster}. First, evidence from Figure~\ref{fig:toy_ae} suggests that shifting the orientation of neurons is the most influential adjustment for modifying semantic representations, aligning directly with OFT’s operational intent. Additionally, OFT’s approach can be interpreted as optimizing neural weights along a smooth hyperspherical manifold, facilitating an optimization landscape that is inherently more favorable. Empirical findings supporting this perspective are provided in \cite{liu2021orthogonal}.

\textbf{Why not minimize hyperspherical energy}. Distinct from \cite{liu2021orthogonal}, our objective does not involve the minimization of hyperspherical energy. For classification applications, it is preferable for neurons to exhibit minimal redundancy. Achieving the lowest possible hyperspherical energy corresponds to an even distribution of neurons on the hypersphere, which, although theoretically appealing, is counterproductive during finetuning as it may erase valuable information gained from pretraining.

\textbf{Balancing flexibility and regularization in finetuning.} Our study identifies a fundamental tension between model flexibility and regularization during finetuning. While traditional finetuning achieves the highest level of flexibility, it often sacrifices stability and is prone to causing model degradation or collapse. By contrast, OFT, through its straightforward design, maintains a favorable compromise by enforcing preservation of neuron pairwise angles. Furthermore, employing a block-diagonal structure serves as a more robust form of regularization over the orthogonal transformation.

\textbf{Inference efficiency maintained.} In contrast to approaches such as ControlNet, our OFT strategy introduces no additional computational cost during inference. At test time, the finetuned weights are efficiently constructed by multiplying the trained orthogonal matrix $\bm{R}$ with the original pretrained weight matrix $\bm{W}^0$, yielding a final weight matrix $\bm{W}=\bm{R}\bm{W}^0$. As a result, inference throughput remains identical to that of the original pretrained model.

\section{Experiments and Results}

\textbf{Experimental configuration}. For empirical validation, we employ Stable Diffusion v1.5~\cite{rombach2022high} as the base pretrained text-to-image framework. In order to ensure comparability, image outputs are randomly sampled for each method. For subject-specific image synthesis, we follow the procedures set by DreamBooth~\cite{ruiz2023dreambooth}. For conditional generation, the experimental setup aligns with ControlNet~\cite{zhang2023adding} and T2I-Adapter~\cite{mou2023t2i}. To facilitate a fair comparison against LoRA, both OFT and COFT are restricted to operate on the same layers as those modified by LoRA. Further experimental details and comprehensive results can be found in Appendix~\ref{app:settings}.

\subsection{Subject-driven Generation}

\textbf{Experimental details}. DreamBooth~\cite{ruiz2023dreambooth} and LoRA~\cite{hulora2022} serve as primary baseline methods. All evaluated approaches adopt the loss function described in the DreamBooth framework. For both DreamBooth and LoRA, the implementations adhere closely to their respective original descriptions, utilizing optimal hyperparameter configurations as recommended. Additional findings and supplementary examples are presented in Appendix~\ref{app:settings},~\ref{app:coft_oft},~\ref{app:more_qualitative},~\ref{app:failure}.

\textbf{Finetuning stability and convergence}. We begin by assessing both the stability and convergence rate during finetuning for DreamBooth, LoRA, OFT, and COFT. The corresponding results are presented in Figure~\ref{fig:teaser} and Figure~\ref{fig:db_convergence}. It is evident that COFT and OFT enable the diffusion model to be finetuned in a stable manner. Within 400 training steps, DreamBooth and the OFT-based methods demonstrate effective control over the generation process, yet LoRA struggles to retain subject identity. Progressing to 2000 iterations, DreamBooth tends to produce images with collapsed features, while LoRA is unable to generate the target yellow shirt—substituting yellow fur instead. In contrast, OFT and COFT maintain robust and consistent control over image generation throughout. These observations support that the OFT framework offers not only rapid, but also reliable, convergence for subject-driven generation. A noteworthy advantage is that the stability with respect to the number of finetuning iterations significantly reduces the need for extensive hyperparameter tuning for each subject. This means that for both OFT and COFT, one can confidently assign a large iteration count without precise adjustment. For COFT, as long as an appropriate $\epsilon$ is chosen, neither the learning rate nor the number of iterations require meticulous tuning.

\textbf{Quantitative comparison}. Adopting the methodology from \cite{ruiz2023dreambooth}, we perform quantitative evaluations to measure subject fidelity (using DINO~\cite{caron2021emerging} and CLIP-I~\cite{radford2021learning}), alignment with textual prompts (using CLIP-T~\cite{radford2021learning}), and output diversity (employing LPIPS~\cite{zhang2018unreasonable}). The CLIP-I metric involves computing the mean pairwise cosine similarity between CLIP embeddings of real and generated images. For DINO, a similar approach is used, replacing CLIP embeddings with ViT S/16 DINO embeddings. CLIP-T assesses the average cosine similarity between embeddings of text prompts and their corresponding generated images. Sample diversity is quantified by the average LPIPS cosine similarity among images generated from identical prompts and subjects. As reported in Table~\ref{table:dreambooth_final}, COFT and OFT surpass DreamBooth and LoRA by a substantial margin in both DINO and CLIP-I measures, while also delivering equal or slightly superior results for text prompt fidelity and diversity. Each technique is evaluated by finetuning the same pretrained model 30 times using distinct random seeds to ensure results are not affected by stochastic variation. Collectively, these findings indicate that our OFT framework achieves superior stability and convergence while consistently producing enhanced performance.

\textbf{Qualitative comparison}. To provide a clearer insight into the advantages of OFT, we present several randomly selected cases of subject-driven image generation in Figure~\ref{fig:dreambooth_final}. For consistency and fairness, all sample outputs are derived from the same finetuned base model across different techniques, and none of the text prompts are individually tuned to enhance their outputs. For each technique, we select the checkpoint corresponding to the highest validation CLIP score. Inspection of Figure~\ref{fig:dreambooth_final} indicates that both OFT and COFT are highly effective at maintaining semantic fidelity to the subject, whereas LoRA frequently struggles to keep the subject identity intact (\emph{e.g.}, the bowl example where LoRA entirely loses the subject characteristics). Furthermore, both OFT and COFT exhibit superior text prompt controllability compared to DreamBooth, which, while it generally preserves the subject, often does not faithfully follow the prompt (\emph{e.g.}, the bowl example’s first row). These visual comparisons illustrate that OFT excels by simultaneously delivering strong controllability and robust subject preservation. Additionally, the performance of OFT remains stable even when the number of training iterations is increased, unlike DreamBooth and LoRA, which show diminished effectiveness under similar circumstances. Additional visual results can be found in Appendix~\ref{app:more_qualitative}. In Appendix~\ref{app:human}, we also include a human assessment that further corroborates the advantages of OFT.

\subsection{Controllable Generation}

\textbf{Settings}. For baseline comparison, we utilize ControlNet~\cite{zhang2023adding}, T2I-Adapter~\cite{mou2023t2i}, and LoRA~\cite{hulora2022}. The principal paper evaluates three difficult controllable generation scenarios: Canny edge to image~(C2I) on the COCO dataset~\cite{lin2014microsoft}, segmentation map to image~(S2I) on ADE20K~\cite{zhou2017scene}, and landmark to face~(L2F) on CelebA-HQ~\cite{karras2018progressive,xia2021tedigan}. Each method finetunes the Stable Diffusion~(SD) v1.5 model on these tasks for 20 epochs. Further experiments and results are available in Appendix~\ref{app:more_qualitative},\ref{app:more_control_task},\ref{app:failure}.

\textbf{Convergence}. To assess convergence rates, we analyze ControlNet, T2I-Adapter, LoRA, and COFT on the L2F benchmark. Both quantitative and qualitative methods are employed for this analysis. For the quantitative aspect, the main metric is the average $\ell_2$ distance calculated between the control and the generated face landmarks. Figure~\ref{fig:face_convergence} visualizes the trajectory of face landmark error as each model is fine-tuned across a range of epochs. The data indicates that COFT and OFT both demonstrate markedly more rapid convergence compared to the alternatives. For example, LoRA only attains the level of performance OFT reaches at epoch 8 after 20 epochs of training. It is noteworthy that OFT and COFT are designed with a parameter count similar to that of LoRA (and substantially lower than that of ControlNet), but they achieve convergence far more efficiently relative to previous approaches. Additionally, the accelerated convergence of OFT is corroborated in Figure~\ref{fig:teaser}, where the right-side example highlights the superior data efficiency of OFT: it converges effectively using merely 5\% of the ADE20K data, outperforming both ControlNet and LoRA in this regard. In terms of qualitative evaluation, our comparison centers on OFT, COFT, and ControlNet, as ControlNet is the closest competitor in terms of landmark error. As shown in Figure~\ref{fig:face_convergence_qual}, both OFT and COFT exhibit stable convergence behaviors; the generated face poses gradually synchronize with the target control landmarks as training progresses. Conversely, with ControlNet, improved controllability only becomes apparent after epoch 8—a behavior consistent with the abrupt convergence phenomenon detailed in \cite{zhang2023adding}. The stable and gradual convergence dynamics of the OFT framework make it significantly more user-friendly in practice, as its training process is smoother and easier to anticipate. This in turn simplifies the task of tuning OFT hyperparameters.

\textbf{Quantitative comparison}. To assess controllable generation, we propose a metric termed control consistency, which quantifies how well the control signal extracted from the generated output matches the initial input control signal. Specifically, for the C2I task, we employ IoU and F1 score, while for the S2I task, we report both mean IoU and mean as well as overall accuracy. In the case of the L2F task, we measure the average $\ell_2$ distance between the predicted and original control landmarks. Further specifics regarding these metrics can be found in Appendix~\ref{app:settings}. Each baseline and comparison method utilizes its optimal hyperparameter configurations. According to Table~\ref{table:control_gen}, both OFT and COFT considerably outperform other baselines in terms of control strength and precision. Notably, methods based on adapters (\emph{e.g.}, T2I-Adapter and ControlNet) demonstrate slower convergence rates and inferior performance at convergence. LoRA achieves better results than ControlNet on the S2I benchmark, but is surpassed by ControlNet in both C2I and L2F. Overall, the maximum achievable performance with LoRA remains modest, despite exhaustive tuning of its parameters. Conversely, OFT continues to benefit from increases in trainable parameter count, and has yet to reach its optimal potential. Our methodology for the quantitative assessment of controllable generation represents a key contribution, providing an accurate gauge of the control ability in finetuned models, as determined by the precision of the pre-trained segmentation or detection networks.

\textbf{Qualitative comparison}. We further conduct a qualitative assessment of OFT and COFT against leading methods such as ControlNet, T2I-Adapter, and LoRA. As depicted by the randomly sampled outputs in Figure~\ref{fig:control_final}, OFT and COFT consistently produce images that are not only highly realistic and visually convincing, but also exhibit precise controllability. On the S2I task, it is evident that LoRA fails to adhere to the structure dictated by the input segmentation map, whereas ControlNet, OFT, and COFT all successfully impose control over the generated imagery. Compared to ControlNet, OFT and COFT generate images with enhanced visual detail and more coherent geometric structures, while utilizing significantly fewer parameters. In the C2I task, OFT and COFT are capable of synthesizing plausible and semantically aligned images from basic Canny edge inputs, outperforming both T2I-Adapter and LoRA by a substantial margin. For the L2F task, our approach achieves the most accurate reproduction of facial poses, even for complex and challenging viewpoints. Across all three control scenarios, OFT and COFT surpass state-of-the-art baselines in image quality, underscoring the strength of the OFT approach for controlled image synthesis. Additional qualitative samples for each control task are presented in Appendix~\ref{app:control_exp}. Furthermore, Appendix~\ref{app:more_control_task} illustrates OFT's versatility by demonstrating successful application to a broader set of control tasks, including dense pose-to-human body, sketch-to-image, and depth-to-image synthesis.

\section{Concluding Remarks and Open Problems}

This work is inspired by the finding that angular relationships between neurons play a pivotal role in shaping visual semantics. To this end, we introduce a straightforward yet powerful finetuning approach—orthogonal finetuning—for enhancing control over text-to-image diffusion models. Our method is particularly aimed at two key tasks: subject-driven generation and controllable generation. In comparison to prior approaches, OFT achieves superior levels of control and stability during finetuning while reducing the number of tunable parameters required. Crucially, OFT incurs no extra inference cost, enabling the resulting models to be deployed efficiently.

Despite its strengths, OFT opens several avenues for further exploration. The technique enforces orthogonality through Cayley parametrization, which necessitates a matrix inversion. This requirement slightly constrains OFT's scalability. Although we mitigate this with a block diagonal parametrization, developing a faster, differentiable method for matrix inversion remains an open technical hurdle. Additionally, OFT holds promise for compositional finetuning: orthogonal matrices from several OFT finetuning sessions can be multiplied together, with the product retaining orthogonality. Investigating whether this collection of orthogonal matrices effectively consolidates the knowledge from all involved tasks is an intriguing research question. Lastly, OFT's parameter efficiency currently hinges on the block diagonal architecture, which brings about inherent biases and reduces flexibility. Discovering methods to enhance parameter efficiency while minimizing bias is another outstanding challenge.

\end{document}